\section{Related work}
\label{sec:related}

\subsection{Risk-aware autonomous intersection planning}

Chance-constrained planning converts a probabilistic forecast into a bound on
constraint violation, but the interpretation of that bound depends on the
uncertainty model and on whether the mathematical assumptions remain valid in
deployment. Iterative risk allocation was introduced to redistribute a joint
chance budget toward active constraints rather than impose an equally
conservative allocation everywhere \citep{ono2008risk}. In driving, robust
scenario MPC represents uncertain multimodal obstacles with multiple scenarios
\citep{batkovic2021scenario}; Branch MPC instead plans a trajectory tree over
reactive modes and uses a coherent risk measure to trade performance against
robustness \citep{chen2022branch}. Both illustrate that ``risk'' is an
optimisation design variable, not the same quantity as measured collision or
physical clearance.

The closest control lineage to this dissertation is the SMPC formulation of
\citet{nair2022smpc}, which embeds Gaussian-mixture predictions in multimodal
chance constraints and evaluates feedback-policy variants at a CARLA
intersection. Its later extension jointly optimises feedback policies and
mode-wise risk allocation and demonstrates real-time closed-loop operation in
several scenarios \citep{nair2025predictive}. Complementary work establishes
recursive-feasibility and safety properties for nominal, robust and
contingency chance-constrained planners under explicit Gaussian-mixture
assumptions \citep{ren2025recursively}. These are stronger formal and
cross-scenario contributions than the empirical study here; however, a formal
property of the planned policy does not by itself quantify executed
bounding-box separation after solver failure handling and safety supervision.

The present work therefore asks a narrower empirical question. Adaptive risk
is compared with aggressive, medium and conservative fixed operating points,
not with a single convenient baseline. Predictor stack and target style are
crossed with policy, and executed completion, physical separation, binary
failures, solver identity and supervisor activity are retained. This design
addresses RQ4 while avoiding the stronger claim that a finite simulator matrix
validates a chance constraint or proves real-world safety.

\subsection{Multimodal prediction, scoring and adaptation}

MultiPath represents future motion by a discrete distribution over trajectory
anchors with regressed offsets and uncertainty \citep{chai2020multipath}. This
compact Gaussian-mixture output is particularly suitable for an optimisation
consumer. MultiPath++ subsequently replaced dense raster-only fusion with
sparse scene elements, learned anchors and output aggregation
\citep{varadarajan2022multipathpp}. The dissertation retains the original
MultiPath decoder and frozen anchors so that adaptation changes a deployed
component without replacing the planner interface.

Displacement and probabilistic scores answer different questions. Top-1 ADE
and FDE assess one selected trajectory and can ignore whether probability mass
is distributed appropriately across plausible modes. Negative log likelihood
uses the full predictive density and is a proper score under the relevant
distributional conditions \citep{gneiting2007scoring}, but a low value can
arise through both better means and changed covariance. Calibration is again
distinct: neural-network confidence can be systematically miscalibrated
\citep{guo2017calibration}, and post-hoc calibrated regression can improve
coverage without changing the underlying point predictor
\citep{kuleshov2018calibrated}. Subgroup calibration matters because an
aggregate guarantee need not hold on identifiable subpopulations
\citep{hebertjohnson2018multicalibration}; trajectory-specific work likewise
finds that uncertainty estimates require calibration before planning
\citep{cao2024cctr}. These distinctions motivate reporting ADE, FDE,
uncalibrated NLL, coverage and calibrated NLL, with the response-active tail
kept separate from the aggregate.

Adaptation is also a distribution question rather than merely an architecture
choice. UniTraj reports substantial cross-dataset degradation and improvement
from greater data diversity \citep{feng2024unitraj}. Here, the intended claim
is deliberately in-distribution: rollouts are disjoint by initialisation but
come from one Town05 give-way process. Retraining the final prediction head is
therefore a task-adaptation control for domain mismatch, not evidence of
cross-map generalisation.

\subsection{Interaction-aware and Transformer prediction}

Self-attention provides content-dependent interaction across a sequence
\citep{vaswani2017attention}. Driving predictors have used it at substantially
greater scale and with richer representations than this dissertation.
AgentFormer jointly attends across time and agents
\citep{yuan2021agentformer}; Scene Transformer supports marginal, joint and
ego-conditioned multi-agent prediction in a unified architecture
\citep{ngiam2022scene}; AutoBots models socially consistent joint futures as a
sequence of agent sets \citep{girgis2022autobots}. Wayformer studies early,
late and hierarchical fusion of heterogeneous scene modalities
\citep{nayakanti2023wayformer}, while MTR couples global intention queries with
local motion refinement \citep{shi2022mtr}. These results establish attention
as a strong design family, not a default optimum for every data regime.

Architecture comparisons are confounded when representation, residual output,
parameter count or optimisation budget changes simultaneously. The current
study therefore uses two MLP/Transformer pairs with the same input and
residual-output scope and approximately similar adapter capacity. It also
retains a simpler, larger task-adapted head and reports every trainable count
and boundary-epoch selection. Frozen-input zeroing and shuffling test whether
the Transformer actually uses its six-token history, but cannot prove semantic
interaction reasoning. RQ2 is consequently about the tested adapters under a
fixed small-data budget; it is not a referendum on the larger models above.

\subsection{Prediction-aware planning and closed-loop evaluation}

Prediction and planning can be coupled during training or evaluation.
Topological multimodal prediction has been demonstrated with an
optimisation-based navigator at uncontrolled CARLA intersections
\citep{roh2021topological}. DIPP back-propagates through a differentiable
planner and reports both open- and closed-loop gains
\citep{huang2023dipp}. The nuPlan benchmark similarly argues that open-loop
distance metrics are insufficient for long-horizon planning and supplies
reactive closed-loop evaluation \citep{caesar2021nuplan}. These systems cover
broader planning tasks than the controlled junction here, but they also change
more of the stack; they do not isolate how a frozen probabilistic predictor
interacts with a risk-policy frontier.

Task relevance provides the key evaluation principle. \citet{farid2023task}
propagate prediction failures to planning cost rather than treating every
forecast error as equally harmful. \citet{nakamura2025regret} use closed-loop
regret to identify prediction failures that degrade the robot, and find that a
small high-regret subset can be informative for fine-tuning. Most directly,
the recent preprint by \citet{bouzidi2025closing} reports that open-loop model
ranking need not predict closed-loop performance and that smaller predictors
can match larger ones in its tested stack. The present dissertation agrees with
that broad warning, but contributes a different controlled contrast: it first
verifies that B1 remains better than B0 in-loop, then tests translation within
each predictor--risk--target cell and records solver/supervisor mechanisms.

\subsection{Synthesis and study positioning}

Across the reviewed literature, four claims form a hierarchy rather than a
single model leaderboard: adaptation can improve a forecast; an interaction
architecture can add value beyond matched controls; an offline improvement can
benefit closed-loop execution; and adaptive risk can dominate fixed operating
points. Prior work supplies strong methods for each layer, but the directly
comparable studies reviewed above do not evaluate all four with a frozen
predictor interface, rollout-disjoint selection, several fixed-risk controls,
paired target styles and executed mechanism telemetry. This dissertation fills
that bounded empirical gap. Its contribution is a controlled claim audit in a
single give-way system, not a new universal predictor or a formal safety proof.
